\section{群間計画(Between Design)と分散分析}
\label{b18_BetweenDesign}

\subsection{授業内容}
\subsubsection{概要}

前回学んだ実験計画法の基礎知識を踏まえ，本コマでは群間計画(Between-subjects Design)に焦点を当てる。群間計画とは，異なる実験参加者グループに異なる条件を割り当てる実験デザインであり，心理学研究でよく用いられる手法である。特に，一要因および二要因の群間計画における効果の分解方法と評価方法を学び，分散分析(ANOVA)の基本的な考え方を理解する。平方和の分解，交互作用の概念について学習する。

\subsubsection{コマ主題細目}
\begin{description}
	\item[一要因群間計画の平方和の分解] 群間計画(Between-subjects Design)は，異なる実験参加者グループに異なる条件を割り当てる実験デザインである。群間計画では，各参加者は一条件のみを経験するため，練習効果や順序効果といった問題を回避できるが，個人差による誤差が大きくなる傾向がある。分散分析の核心は，データの総変動(平方和)を複数の変動源に分解することにある。一要因群間計画では，総平方和($SS_T$)は要因の効果による平方和($SS_A$)と誤差による平方和($SS_E$)に分解される($SS_T = SS_A + SS_E$)。各平方和の計算方法と意味を理解し，効果と誤差をどのように分離するかを学ぶ。
	      \begin{flushright}
		      $\to$ \textcite{Toyoda201706}のP.27--44，\textcite{kosugi2018}のP.74--78.
	      \end{flushright}

	\item[二要因群間計画と主効果] 二要因群間計画では，一要因の場合に加えて，第二要因($SS_B$)と交互作用($SS_{AB}$)の平方和が加わる($SS_T = SS_A + SS_B + SS_{AB} + SS_E$)。全体平均，群平均，セル平均の関係と，それらを用いた効果の算出方法についても理解する。
	      \begin{flushright}
		      $\to$ \textcite{Yamada2004}のP.184--189，\textcite{kosugi2018}のP.81--88.
	      \end{flushright}


	\item[交互作用の理解と解釈] 二要因以上の実験計画では，要因間の交互作用(interaction)が重要な意味を持つ。交互作用とは，一方の要因の効果が他方の要因の水準によって異なる現象を指す。主効果が周辺平均の比較であるのに対し，交互作用はセルの中での比較であり，いずれも全体平均からの相対的な比較であるから効果の総和はゼロになる。交互作用の種類，交互作用の平方和を算出する手順を学ぶ。交互作用の視覚的理解を助けるグラフの描き方と読み方についても習得する。
	      \begin{flushright}
		      $\to$ \textcite{Yamada2004}のP.190--193，\textcite{kosugi2019}のP.139--142.
	      \end{flushright}

	\item[平方和の分解と変動の源泉] 要因の数が増えると，平方和の分解に現れる交互作用の組み合わせが急速に増える。3要因になると，3つの主効果に加え，AB，AC，BCの一次の交互作用だけでなく，ABCの二次の交互作用も考える必要がある。要因計画はいくらでも複雑にしてよいものではなく，要因が増えるごとに考えるべき要素が増え，統一的な解釈が難しくなることを理解する。


\end{description}
\subsubsection{キーワード}
\begin{itemize}
	\item 群間計画(Between-subjects Design)
	\item 平方和の分解と分散分析
	\item 主効果と交互作用
\end{itemize}
\subsection{授業情報}
\paragraph{コマの展開方法}
講義
\paragraph{公認心理師標準シラバスにおける位置づけ}
科目番号5 心理学統計法；①心理学で用いられる統計手法；H 推測統計：代表値と散布度をめぐって(2) に相当
\paragraph{心理統計法標準カリキュラム(日本心理学会)における位置づけ}
第6回 実験的研究と要因計画 細目2，第11回 分散分析 細目1・3 に相当
\subsubsection{予習・復習課題}
\paragraph{予習}
前回学んだ実験計画法の基礎概念(要因と水準，ランダム化など)を復習しておくこと。また，線形モデルの基本的な考え方(切片と傾き，回帰式)についても再確認しておくと理解が深まる。さらに，分散(variance)の概念と計算方法について復習しておくことも有用である。

\paragraph{復習}
授業で学んだ内容を定着させるために，以下の点について理解を深めておくこと：

\begin{enumerate}
\item 一要因および二要因の群間計画における平方和の分解方法を理解し，具体例を用いて計算できるようにする
\item 交互作用の概念を理解し，様々なパターンの交互作用(相乗効果，拮抗効果，クロスオーバーなど)を図示して説明できるようにする
\item 提示された実験デザインを適切に表記し，その構造(要因，水準，群間/群内)を説明できるようにする
\end{enumerate}

次回の授業では，Rを用いた群間計画データの分析を行うため，統計環境Rの基本操作を再確認しておくことが望ましい。特に，データフレームの操作方法や基本的な関数の使い方について復習しておくと良い。